% ARGUES: the mechanistic half of the thesis chain.
% DRAWS FROM: PRE_REGISTRATION_OSF.md §2 (H1.3, H1.4), §5 (procedure),
%             §6 (E3 guardrails, E4 model-choice and confound rules).
% GATED ON: implementation. Scope decided 2026-08-12: both get built, so the §2
%           single-publication commitment holds and no scope deviation is declared.
%
% Build order and why: E4 first, because it trains on the Phase C generations
% that already exist and needs no new model inference. E3 second, because it
% needs activation capture and must clear the pre-registered probe sanity gate
% before it touches any citation data. E3 ran 2026-09-05/06 on the registered
% 8B scale-check arm (4-bit, local GPU); the 70B primary is uncollected.
% Numbers: data/analysis/phase_c/e3_results.json (src/e3/run_e3.py).

\section{Mechanistic experiments}

\subsection{E3: internal-state probe}

\guide{H1.3: a linear probe on hidden states at the answer-token positions
separates real from fabricated citations at AUROC of at least 0.70 on
Tülu-3-70B-Instruct.

The null interpretation is pre-registered and unusually specific, so it cannot
be softened after the fact: a null is a cross-task probe-transfer failure or
non-linear decodability, not evidence that no internal fabrication signal
exists. Wen et al. is the precedent for that reading and the footnote scoping
the parallel to probe non-transfer is required.

Pre-declared null title: \emph{Citation-Fabrication Signal Not Linearly
Decodable by a General-Truthfulness Probe: A Cross-Task Transfer Limitation}.

The guardrails in pre-reg §6 on probe decodability versus knowledge apply to
every sentence in this subsection.}

\label{sec:e3}

E3 was run on the registered scale-check arm, Tülu-3-8B, and not on the
registered primary, Tülu-3-70B-Instruct. The reason is compute: the study's
remaining budget was zero dollars, the local GPU holds 12~GB, and a 70B model
does not fit it at any precision that preserves the hidden states. The 8B
checkpoint was loaded in 4-bit NF4 quantization for the same reason, so the
activations read here come from a quantized copy of the weights that generated
the Phase C answers, not the bf16 weights themselves. Both facts are recorded in
Appendix~\ref{app:deviations}. H1.3 as registered, on the 70B primary, is
therefore untested; what follows is the registered procedure applied to the
registered scale-check model.

\paragraph{Procedure.}
The probe is a logistic regression on standardized mean-pooled hidden states.
It was trained on 4{,}026 true/false statements built from TruthfulQA and SciQ,
one true and one false statement per source question, with 502 held out for
validation; no probe statement overlaps any citation-elicitation question. The
layer was chosen on that validation split alone, never on citation data: layer
14 of 33, where the probe separates true from false statements at AUROC
$0.947$. The registered sanity gate asks the probe to reproduce the published
performance of \citet{orgad2024} before it touches citations. They report
exact-answer-token probe AUCs of $0.83$ to $0.95$ across TriviaQA, Winobias and
Math on 7B and 8B instruction-tuned models; our $0.947$ is on a different
benchmark and a different model, so it matches their published range rather
than reproducing any one figure, and the gate is passed in that sense only.
The citation set is
every verified citation in Tülu-3-8B's own Phase C generations whose raw span
could be located in the answer text: 4{,}587 citations, 2{,}363 fabricated and
2{,}224 real. Span location succeeded for $95.49\%$ of real and $94.14\%$ of
fabricated citations, so the two classes are lost at nearly the same rate. The
230 UNVERIFIABLE citations are in neither class and were excluded. No Phase C
citation carries the MISGROUNDED label, so the registered ``fabricated plus
misgrounded'' class is FABRICATED alone. Each citation's hidden states are the
layer-14 states mean-pooled over its own tokens, read in the exact prompt the
generation was produced under.

\paragraph{The cross-task probe does not transfer.}
Applied to the citation activations, the truthfulness probe scores AUROC
$0.434$. This is not a layer-choice artifact: across all 33 layers the citation
AUROC never exceeds $0.528$, while the same probes climb from $0.58$ to $0.95$ on
their own validation split, so the probe learns its own task and transfers
nowhere. Twenty label-shuffled refits, the registered leak control, span
$0.398$ to $0.625$ with mean $0.500$, and 30\% of them sit at least as far from
$0.5$ as the observed $0.434$. The cross-task probe is indistinguishable from a
probe that never saw a label. This is the pre-declared null, and it carries the
pre-declared reading: a cross-task transfer failure of a general-truthfulness
probe, not evidence that no fabrication signal is represented.\footnote{The
parallel to \citet{wen2024} is scoped to probe non-transfer: their probe for
intended sophistry did not generalize to the unintended case. It is not a claim
that the internal signal is absent.}

\paragraph{The within-domain probe does.}
The registered secondary check trains the same probe class on citation
activations and evaluates it on citations from disjoint questions, split by
question so no question straddles the boundary (2{,}294 training and 2{,}293
test citations). It scores AUROC $0.886$. The pre-committed agreement rule then
fires: the within-domain probe exceeds the cross-task probe by $0.453$, far past
the $0.10$ margin, so the cross-task result is downgraded to ``may be partly
reading task format'' and the within-domain probe becomes the headline E3
estimate. Against the registered threshold of $0.70$, \textbf{H1.3 is met by the
headline estimate on the scale-check arm}.

\paragraph{Difficulty deconfound.}
The registered guardrail requires the fabrication probe to out-perform a probe
trained on a per-question difficulty proxy before the signal may be called
fabrication rather than difficulty. The proxy is the registered one, the base
model's log-probability of producing any citation for the item, estimated from
Llama-3.1-8B's own Phase C generations: four draws per question, one per prompt
condition, passed through the same extractor as everything else, smoothed as
$\log((k + \tfrac12)/(n + 1))$. Four draws give five distinct levels, which is
coarse for a quantity the registration calls continuous, and we say so rather
than dress it up. A ridge regression from the same layer-14 states to the proxy,
scored as a fabricated-versus-real classifier on the same held-out citations in
whichever orientation is larger, reaches $0.522$; the raw proxy with no hidden
states at all reaches $0.504$. On the registered high-difficulty cells,
obscure-but-real and unanswerable (1{,}461 held-out citations), the fabrication
probe scores $0.895$ against $0.530$ for the difficulty probe and $0.572$ for the
raw proxy. Item difficulty as proxied here barely predicts which citations are
fabricated, and the hidden states carry essentially none of it, so the
within-domain signal is reported as fabrication-decodability.

\paragraph{What this licenses.}
On Tülu-3-8B, a fabrication-correlated signal is linearly decodable at the
answer tokens from disjoint questions, and it is not item difficulty. That is
the whole claim. It does not say the model represents or knows that a citation
is fake; the guardrails in pre-registration §6 forbid that sentence, and the
design cannot support it. It does not transfer from general truthfulness, which
is the registered primary probe and the registered null. And it is a result on
the 8B scale-check arm read from quantized weights, with the 70B primary and the
Llama-3.1-70B-Instruct cross-pipeline check both uncollected.

\subsection{E4: naive-judge cross-family}

\guide{H1.4: a citation-reviewer preference model trained only on surface
features (count, format, plausible author and venue density, explicitly not
existence) prefers post-trained over base outputs at least 60\% of the time on
cross-family transfer, training on Tülu-3 and evaluating on OLMo-2.

This is the plausibility argument for the reward proxy, and the pre-registration
is emphatic that it is not causal identification. Keep that label on it.

Pre-declared null title: \emph{Naive Citation-Quality Proxies Do Not Prefer
Post-Trained Outputs: A Null on the Plausible-Proxy Hypothesis for Citations}.}

The preference model was trained on 14{,}400 Tülu-3 generations, 7{,}200 of them
post-trained, and evaluated on 3{,}600 held-out OLMo-2 base-versus-Instruct
pairs. Training and evaluation families are disjoint, so every number below is
cross-family transfer.

\paragraph{The registered test.}
The four-feature hand-coded logistic prefers the post-trained output in
$85.78\%$ of pairs, bootstrap CI $[85.00, 86.61]$. The DeBERTa-v3-base arm
reaches $87.24\%$, $[86.19, 88.29]$. DeBERTa leads by $1.46$ points, inside the
pre-registered 5-point band, so the model-choice rule of §5 makes the
\emph{logistic} the primary reported model. \textbf{H1.4 is supported}: the rate
clears the registered 60\% threshold and the CI lower bound of $85.00\%$ sits far
above chance. Under the registered count-matched ($\pm 1$) control the logistic
falls to $78.76\%$, $[77.66, 79.81]$, whose lower bound lies between 50\% and the
unconstrained lower bound. By the registered rule that is a
\emph{count-sensitive} plausible-proxy preference: the citation-count confound is
weakened, not excluded.

\paragraph{Both registered controls turned out weaker than their names suggest.}
Neither problem is visible from the registered numbers alone, and both were
found by inspecting the control rather than the result. First, the $\pm 1$ window
is permissive on this corpus. Base checkpoints average $0.09$ citations per
generation and final checkpoints $1.31$, so a $\pm 1$ window still admits the
0-against-1 comparison, which is the dominant contrast in the data. The control
removes the large count gaps and leaves precisely the one carrying the effect.
Second, the placebo is not direction-free as specified. It pairs SFT against DPO,
which differ in stage, and citation surface is still rising across the cascade
(citation count $0.758 \to 1.329$, format conformity $0.454 \to 0.636$,
plausible-author density $0.291 \to 0.407$, plausible-venue density
$0.385 \to 0.528$). Its observed rate of $64.71\%$ is therefore ambiguous between
a spurious directional bias in the judge and the judge correctly reading a real
increase. As specified, it cannot separate the two.

\paragraph{Holding citation count exactly equal resolves both.}
Restricting to pairs whose citation counts are identical removes the count
channel entirely, so any residual preference must come from the other three
features. This tier is post-hoc and is labeled as such in the artifact; the
registered control remains the $\pm 1$ window reported above.

\begin{table}[htbp]
\centering
\small
\caption{E4 preference rate by count-matching tier, logistic arm. The registered
control is the $\pm 1$ tier; the exact-match tier is a post-hoc diagnostic.}
\label{tab:e4count}
\begin{tabular}{@{}l l l@{}}
\toprule
Tier & Primary (base vs Instruct) & Placebo (SFT vs DPO) \\
\midrule
All pairs                  & 0.8578 $[0.8500, 0.8661]$, $n = 3600$ & 0.6471 $[0.6347, 0.6589]$, $n = 3600$ \\
Within $\pm 1$ (registered) & 0.7876 $[0.7766, 0.7981]$, $n = 2283$ & 0.5808 $[0.5683, 0.5932]$, $n = 2914$ \\
\textbf{Exactly equal}     & \textbf{0.5155} $[0.5083, 0.5238]$, $n = 904$ & \textbf{0.5035} $[0.4923, 0.5146]$, $n = 1882$ \\
\bottomrule
\end{tabular}
\end{table}

Two readings follow, and they point the same way. The placebo comes out clean
once counts are equal: $50.35\%$ with a CI straddling $0.50$. The judge carries
no directional bias of its own, and its entire placebo signal was the real
SFT-to-DPO citation increase. The control passes once it is asked a question it
can answer. The primary preference, meanwhile, nearly vanishes: $51.55\%$ on 904
pairs of which 856 are ties, leaving only 48 pairs decided at all, 38 for the
post-trained side and 10 for base.

\paragraph{What this licenses, and what it does not.}
The naive judge's preference for post-trained output is almost entirely
\emph{cites more}, not \emph{cites in a more plausible-looking way}. Format
conformity, author-name plausibility and venue plausibility add close to nothing
once quantity is held fixed. This does not weaken E4's role in the argument: E4
exists to show that a plausible reward proxy would reward the post-trained
behavior, and citation quantity is an entirely ordinary thing for a reward model
to reward. A proxy keying on ``did it cite sources'' prefers post-trained output
at 86\%, and that is the mechanism the paper points at. It does change what may
be claimed. The stronger reading, that a naive judge is fooled by
fabricated-looking scholarly apparatus, is \textbf{not supported} here, and we do
not claim it. E4 supports ``quantity is rewarded,'' not ``surface plausibility is
rewarded.''

Two limitations belong with this result. The exact-count tier is post-hoc, added
after the $\pm 1$ window proved permissive on this data; the registered control
is the $\pm 1$ result and it stands as the registered one. And the DeBERTa arm
has no exact-match number, because the fine-tuned weights were not persisted and
retraining would produce a second version of a one-shot registered quantity.
Since the logistic is the primary model, that gap sits in the secondary check
only.
